\chapter{Executive Summary}

\section{Overview}
This report presents the results of the penetration testing engagement conducted by the student \textbf{Mariano Aponte}, with student ID \textbf{522501877} for the \textbf{Penetration Testing and Ethical Hacking exam}. 
\\

The assessment was carried out between \textbf{August 1, 2026} and \textbf{August 30, 2026}, and targeted the \textbf{"Corrosion: 1"} asset, a deliberately vulnerable system available on the VulnHub platform \cite{vulnhubCorrosion}.
\\

The objective of the assessment was to evaluate the security posture of the asset by identifying vulnerabilities that could be exploited by a malicious actor.
\\

The assessment was performed using a combination of automated and manual testing techniques, in a black-box fashion. The methodology was aligned with the \textbf{Penetration Testing Execution Standard (PTES)} \cite{ptes}. 
\\

All the testing activities were performed in accordance to the Rules of Engagement agreed by both parties, available in the Engagement Highlights section.
\\

The information available in this document is meant to be interpreted as educational only. All the penetration testing activities were carried out for educational purposes. No third parties assets were harmed or tested without explicit authorization.

\section{Overall Assessment}

The assessment identified \textbf{89 vulnerabilities} of the following severity:
\\


\begin{center}
	\begin{tabular}{lr}
		\toprule
		\textbf{Severity} & \textbf{Count} \\
		\midrule
		Critical & 21 \\
		High & 45 \\
		Medium & 19 \\
		Low & 0 \\
		Not Rated/Info & 4 \\
		\bottomrule
	\end{tabular}
\end{center}

Several of these vulnerabilities were successfully chained during exploitation, ultimately resulting in complete compromise of the target system.
\\

Based on the number and severity of identified vulnerabilities and the successful compromise of the host, the overall security posture is assessed as \textbf{Unacceptable}.

\section{Key Findings}

The most significant findings identified during the assessment include:

\begin{enumerate}
\item \textbf{Critical: Web Application Vulnerability Leading to Remote Access}

A vulnerability in the web application allows unauthorized access to sensitive files and can be subsequently used to execute commands on the server. This weakness could provide an attacker with \textbf{remote access to the system}.

\item \textbf{Critical: Outdated Apache HTTP Server with Known Vulnerabilities}

The target runs \texttt{Apache HTTP Server 2.4.46}, an outdated version affected by multiple known vulnerabilities. Depending on the enabled modules and configuration, these vulnerabilities could lead to denial of service and/or remote code execution.

\item \textbf{Critical: Local Privilege Escalation}

A user utility has excessive privileges and could be manipulated to execute unauthorized commands with \textbf{root-level privileges}. This weakness could allow an attacker to escalate from a standard user account to full administrative access.

\item \textbf{High: Exposure of Sensitive Files and Authentication Information}

Several files are accessible through the web application and can disclose information useful for further attacks. In particular, the application exposes \texttt{/var/log/auth.log}, which contains authentication activity, access attempts, and other security-relevant information.

\item \textbf{High: Weak Protection of Sensitive Backup Data and Credentials}

A password-protected backup archive containing user credentials is discovered on the system. The archive password is weak and can be easily cracked, allowing access to the stored credentials and subsequent authentication as the \texttt{randy} user.

\item \textbf{High: Vulnerable Linux Kernel / Missing Security Patches}

Manual analysis identified a Linux Kernel vulnerability that could allow a local user to gain higher privileges.
\end{enumerate}

Complete technical details and remediation recommendations are provided in the Vulnerability Report section.

\section{Recommendations}

To improve the overall security posture, the following actions are recommended:

\begin{enumerate}
    \item Apply vendor security updates and patches and keep the Operating System and the Apache2 Web Server up to date.
    \item Implement Secure Programming best practices for web-exposed and system applications.
    \item Strengthen password policy for both user accounts and protected files.
    \item Review privilege assignments for critical and high-risk executables and files.
    \item Enhance logging, monitoring, and incident detection capabilities.
    \item Perform regular vulnerability assessments and periodic penetration tests.
\end{enumerate}

\section{Conclusion}
\textbf{This report reflects the security posture of the assessed systems at the time of testing. Changes made to the environment after the assessment may affect the validity of these results. Continuous security monitoring and periodic reassessment are therefore always recommended.}